% ID: SEL-D0081-OA-E01
% Tanda S081.
Calculamos primero los productos y la diferencia:
\[
 A^2=\begin{pmatrix}1&3\\0&4\end{pmatrix},\qquad
 BC=\begin{pmatrix}1&0\\8&1\end{pmatrix},\qquad
 A-BC=\begin{pmatrix}0&1\\-8&1\end{pmatrix}.
\]
Como $\det(A^2)=4$, la ecuación tiene solución única. Multiplicamos por la inversa a la izquierda:
\[
 X=(A^2)^{-1}(A-BC)
 =\begin{pmatrix}1&-3/4\\0&1/4\end{pmatrix}
  \begin{pmatrix}0&1\\-8&1\end{pmatrix}
 =\boxed{\begin{pmatrix}6&1/4\\-2&1/4\end{pmatrix}}.
\]
Comprobamos que la primera fila de $A^2X$ es $(6,1/4)+3(-2,1/4)=(0,1)$ y la segunda es $4(-2,1/4)=(-8,1)$, como exige $A-BC$.
